\input{../tex-inputs/latex-leseansicht-vorspann}

\section[Arthur Schnitzler, Elisabeth Steinruck, Olga Schnitzler und Heinrich Mann an Gustav Schwarzkopf, 7. 9. {[}1912?{]}]{L04456 Arthur Schnitzler, Elisabeth Steinruck, Olga Schnitzler und Heinrich
               Mann an Gustav Schwarzkopf, 7. 9. {[}1912?{]}}
\nopagebreak\mylabel{L04456v}
\rehead{ }\normalsize\beginnumbering\briefempfaengerindex{Schwarzkopf, Gustav@\textsc{Schwarzkopf, Gustav}!zzzMann, Heinrich@\emph{von Heinrich Mann}!1912-09-071@{7. 9. {[}1912?{]}}|(be}\briefempfaengerindex{Schwarzkopf, Gustav@\textsc{Schwarzkopf, Gustav}!zzzSchnitzler, Olga@\emph{von Olga Schnitzler}!1912-09-071@{7. 9. {[}1912?{]}}|(be}\briefempfaengerindex{Schwarzkopf, Gustav@\textsc{Schwarzkopf, Gustav}!zzzSteinrück, Elisabeth@\emph{von Elisabeth Steinrück}!1912-09-071@{7. 9. {[}1912?{]}}|(be}\briefempfaengerindex{Schwarzkopf, Gustav@\textsc{Schwarzkopf, Gustav}!zzzSchnitzler, Arthur@\emph{von Arthur Schnitzler}!1912-09-071@{7. 9. {[}1912?{]}}|(be}
\toendnotes[C]{\smallbreak\pagebreak[2]}
\correspDesc{Versand  durch Arthur Schnitzler, Elisabeth Steinrück, Olga Schnitzler, Heinrich Mann am 7. 9. [1912?] in Tutzing
\newline{}Übermittlung  am 7. 9. [1912?] in Tutzing
\newline{}Erhalt  durch Gustav Schwarzkopf im Zeitraum [8. 9. 1912
                  – 12. 9. 1912?] in Wien}\toendnotes[C]{\smallbreak}
\Standort{DLA, A:Schnitzler, HS.1985.1.1897.}
\physDesc{Bildpostkarte, 285 Zeichen
\newline{}Handschrift Arthur Schnitzler: Bleistift, deutsche Kurrent
\newline{}Handschrift Elisabeth Steinrück: Bleistift, lateinische Kurrent
\newline{}Handschrift Olga Schnitzler: Bleistift, lateinische Kurrent
\newline{}Handschrift Heinrich Mann: Bleistift, lateinische Kurrent
\newline{}Versand: Stempel: »\nobreak{}\oindex{Tutzing@\textbf{Tutzing}|pwk}Tutzing, 7 Sep.\nobreak{}«.  }\toendnotes[C]{\smallbreak}\pstart{}{\pb}\textsc{Herrn Gustav Schwarzkopf}\pend{}\pstart{}Wien I\oindex{I., Innere Stadt@\textbf{I., Innere Stadt}|pw}\pend{}\pstart{}\textsc{Tiefer Graben 17}\oindex{Wien@\textbf{Wien}!I., Innere Stadt@\textbf{I., Innere Stadt}!Tiefer Graben 17@\textbf{Tiefer Graben 17}|pw}\pend{}{\bigskip}
\pstart
           \noindent{}{\pb}\textcolor{gray}{\textbf{Starnberger See\oindex{Starnberger See@\textbf{Starnberger See}|pw}: Tutzing mit Hôtel Simson\oindex{Hotel Simson@\textbf{Hotel Simson}|pw}}}\pend
           \vspace{1em}
\pstart
           \noindent{}{\pb}Herzliche Grüße u auf ſehr bal\textcolor{gray}{dgs}
               Wiederſehn. \label{K_L04456-1v}\edtext{Heut fahren wir}{\lemma{\textnormal{\emph{Heut fahren wir}}}\Cendnote{\textnormal{Dass das
                  Datum des Poststempels auch der Tag des Abfassung der Postkarte war, bestätigt der
                     \emph{Tagebuch}\pwindex{Schnitzler, Arthur 15.\,5.\,1862 Wien – 21.\,10.\,1931 ebd.@\textsc{Schnitzler, Arthur} (15.\,5.\,1862 Wien – 21.\,10.\,1931 ebd.), \emph{Schriftsteller, Mediziner}!Tagebuch@\strich\emph{Tagebuch}|pwk}eintrag vom letzten Tag in Tutzing\oindex{Tutzing@\textbf{Tutzing}|pwk}, in dem die Personenkonstellation und
                  die Abreise dokumentiert sind: »Nm. zu Liesl\pwindex{Steinrück, Elisabeth 19.\,11.\,1885 – 7.\,4.\,1920 Partenkirchen@\textsc{Steinrück, Elisabeth} (19.\,11.\,1885 – 7.\,4.\,1920 Partenkirchen)|pw}; mit Mann\pwindex{Mann, Heinrich 27.\,3.\,1871 Lübeck – 11.\,3.\,1950 Santa Monica@\textsc{Mann, Heinrich} (27.\,3.\,1871 Lübeck – 11.\,3.\,1950 Santa Monica)|pw}.
                     Abschied.«, siehe A. S.: \emph{Tagebuch}, 7. 9. 1912.}}}\label{K_L04456-1} nach Wien\oindex{Wien@\textbf{Wien}|pw}.\pend
           \pstart Ihr \spacefill\mbox{Arthur}\pend{}\selectlanguage{ngerman}\vspace{1em}
\pstart
           \noindent{}{[}hs. Steinrück:{]} Lieber Herr Schwarzkopf, wie schade, dass ich Sie nie seh, aber bald
               komm ich nach Wien\oindex{Wien@\textbf{Wien}|pw}. Herzlichst{\\}Ihre
                  \spacefill\mbox{Lisl St.}\pend
           \selectlanguage{ngerman}\vspace{1em}
\pstart
           \noindent{}{[}hs. O. Schnitzler:{]} Auf sehr baldiges Wiedersehen!\pend
           
\pstart
           Ihre \spacefill\mbox{OlgaS.}{ }{\\[\baselineskip]}\spacefill\mbox{{[}hs. Mann:{]} Heinrich Mann}\pend
           \leftskip=0em{}\selectlanguage{ngerman}\endnumbering\briefempfaengerindex{Schwarzkopf, Gustav@\textsc{Schwarzkopf, Gustav}!zzzMann, Heinrich@\emph{von Heinrich Mann}!1912-09-071@{7. 9. {[}1912?{]}}|)be}\briefempfaengerindex{Schwarzkopf, Gustav@\textsc{Schwarzkopf, Gustav}!zzzSchnitzler, Olga@\emph{von Olga Schnitzler}!1912-09-071@{7. 9. {[}1912?{]}}|)be}\briefempfaengerindex{Schwarzkopf, Gustav@\textsc{Schwarzkopf, Gustav}!zzzSteinrück, Elisabeth@\emph{von Elisabeth Steinrück}!1912-09-071@{7. 9. {[}1912?{]}}|)be}\briefempfaengerindex{Schwarzkopf, Gustav@\textsc{Schwarzkopf, Gustav}!zzzSchnitzler, Arthur@\emph{von Arthur Schnitzler}!1912-09-071@{7. 9. {[}1912?{]}}|)be}\mylabel{L04456h}
\begin{anhang}
\end{anhang}\newcommand{\dateiname}{L04456}\newcommand{\titel}{Arthur Schnitzler, Elisabeth Steinruck, Olga Schnitzler und Heinrich Mann an Gustav Schwarzkopf, 7. 9. [1912?]}\newcommand{\editorInnen}{Herausgegeben von Jahnke, SelmaMüller, Martin Anton}\input{../tex-inputs/latex-leseansicht-abspann}
